\documentclass[11pt]{article}
\usepackage[a4paper,margin=1in]{geometry}
\usepackage{amsmath,amssymb,mathtools}
\usepackage{enumitem}
\usepackage{listings}
\usepackage{xcolor}
\usepackage{booktabs}
\usepackage{array}
\usepackage{longtable}
\usepackage{hyperref}
\hypersetup{colorlinks=true, linkcolor=blue!50!black, urlcolor=blue!50!black,
            bookmarksopen=true, bookmarksnumbered=true}

\definecolor{cbg}{rgb}{0.97,0.97,0.97}
\definecolor{ckw}{rgb}{0.10,0.30,0.70}
\definecolor{cst}{rgb}{0.20,0.50,0.20}
\definecolor{ccm}{rgb}{0.40,0.40,0.40}
\lstdefinestyle{py}{
  language=Python, basicstyle=\ttfamily\footnotesize,
  keywordstyle=\color{ckw}\bfseries, stringstyle=\color{cst},
  commentstyle=\color{ccm}\itshape,
  backgroundcolor=\color{cbg}, showstringspaces=false,
  breaklines=true, breakatwhitespace=false,
  postbreak=\mbox{\textcolor{ccm}{$\hookrightarrow$}\space},
  frame=single, framerule=0pt,
  numbers=left, numberstyle=\tiny\color{ccm}, numbersep=8pt,
  xleftmargin=2.2em, aboveskip=4pt, belowskip=4pt,
}
\lstdefinestyle{sql}{
  language=SQL, basicstyle=\ttfamily\footnotesize,
  keywordstyle=\color{ckw}\bfseries, stringstyle=\color{cst},
  commentstyle=\color{ccm}\itshape,
  backgroundcolor=\color{cbg}, showstringspaces=false,
  breaklines=true, frame=single, framerule=0pt,
  xleftmargin=1.2em, aboveskip=4pt, belowskip=4pt,
}
\lstset{style=py}

\newcommand{\file}[1]{\texttt{#1}}
\newcommand{\code}[1]{\texttt{#1}}

\title{The polilabs knowledge graph, from first principles\\
\large How it evolved, how every layer works, and why --- no condensation}
\author{For Andrew --- written June 12, 2026}
\date{}

\begin{document}
\maketitle
\tableofcontents
\bigskip

\noindent This document builds the polilabs knowledge representation up from
nothing: what problem it solves, the primitives it is made of, every data
source it draws on, every storage layer, the full evolution from the v0.1
schema sketch to the current two-layer architecture, every agent-facing tool,
and walks through the load-bearing code. Part~VII is a cheat sheet indexed by
the questions you have actually asked. Read it linearly the first time; use
Part~VII afterward.

\part{Primitives: the concepts everything else uses}

\section{The problem statement, precisely}
\label{sec:problem}

polilabs exists because of one empirical fact: commercial legal-AI tools
hallucinate on 17--34\% of queries (Magesh et al., Stanford RegLab 2024), and
they do so not because LLMs are hopeless but because the substrate they are
given --- a keyword search box over editorial documents --- is hostile to
programmatic reasoning. The polilabs wager is that if you redesign the
\emph{substrate} (typed entities, verbatim text, explicit provenance, honest
gaps), a large fraction of hallucinations collapse into either correct answers
or honest ``the corpus does not contain this.''

Everything in this document is downstream of three design commitments:

\begin{enumerate}[itemsep=2pt]
  \item \textbf{Verbatim or nothing.} The system never paraphrases law before
        the user sees it. Extraction is mechanical (XML parsing, regex over
        action strings), and anything inferred carries provenance and
        confidence.
  \item \textbf{Honest unknowns are a return type.} ``Not in the corpus,''
        ``no roll call exists,'' and ``status-only record'' are first-class
        answers, not failures.
  \item \textbf{The agent's latency budget is round trips, not milliseconds.}
        See \S\ref{sec:latency}.
\end{enumerate}

\section{Property graphs vs.\ relational tables --- and why polilabs uses both}
\label{sec:graphvsrel}

A \emph{relational table} stores rows with fixed columns; you ask questions by
joining tables on key equality. A \emph{property graph} stores typed nodes and
typed directed edges, each carrying key--value properties; you ask questions
by traversing edges (``from this Bill, follow \code{HAS\_VERSION} then
\code{HAS\_SECTION}'').

The folk claim is that graphs are ``better for connected data.'' The honest
version: graphs win when the question shape is a \emph{path of unknown or
variable length} (``what does this section cite, and what cites that?''), and
relations win when the question shape is a \emph{filtered scan or aggregate}
(``count bills by outcome''). polilabs therefore runs both, over the same
source files:

\begin{center}
\begin{tabular}{@{}lll@{}}
\toprule
 & SQLite (\file{data/polilabs.db}) & K\`uzu (\file{data/polilabs.kuzu}) \\
\midrule
Question shape & filter / aggregate / full-text & traversal / typed neighborhood \\
Examples & counts, outcome sets, name lookup & citation graph, defined-term scoping \\
Latency & $\sim$1--20\,ms & $\sim$10--100\,ms \\
Rebuild & minutes & hours (see \S\ref{sec:incremental}) \\
\bottomrule
\end{tabular}
\end{center}

Neither database is the source of truth. The \emph{files} are
(\S\ref{sec:layers}); both stores are disposable indexes rebuilt from them.
This is the single most important architectural fact about polilabs: you can
delete both databases at any time and lose nothing.

\section{The three retrieval primitives: FTS, embeddings, alias tables}
\label{sec:retrieval}

There are exactly three ways polilabs turns a user's words into rows, and they
have sharply different jobs.

\textbf{Full-text search (FTS5 / BM25).} SQLite's FTS5 builds an inverted
index: for every word (after \emph{Porter stemming}, which maps
``regulating''/``regulates''/``regulation'' to one stem), a list of the rows
containing it. BM25 ranks hits by term rarity and document length. This is the
workhorse for topical queries (``facial recognition in schools'') and costs
$\sim$1--10\,ms.

\textbf{Dense embeddings.} Each section of bill text is mapped by a small
neural model (bge-small-en-v1.5, 384 dimensions) to a vector; a query is
mapped to the same space; cosine similarity finds passages that are
\emph{semantically} close even with zero word overlap (``self-driving cars''
$\to$ ``autonomous vehicles''). polilabs fuses BM25 and dense rankings with
reciprocal-rank fusion (RRF). Embeddings are the only layer with partial
coverage right now (the embed pass is handed off; the coverage report says
exactly how partial --- \S\ref{sec:coverage}).

\textbf{Alias tables.} A user who says ``the CHIPS Act'' is not doing topical
search --- they are naming an entity. The correct primitive is a
\emph{dictionary}: a table mapping every normalized known name to a canonical
ID. The entity-linking literature is unambiguous that dictionary lookup comes
first and covers most mentions (Wikidata alias tables alone resolve
$>$50\% of mentions before any ML), with fuzzy/embedding matching only as a
fallback ranker. polilabs follows exactly this ladder
(\S\ref{sec:lookupcode}).

The design rule that falls out: \emph{name-shaped queries hit the alias
table, topic-shaped queries hit FTS(+dense), and the agent's system prompt
routes between them} (\S\ref{sec:routing}).

\section{Why round trips, not milliseconds, are the latency budget}
\label{sec:latency}

Every query primitive above answers in milliseconds. But the agent consuming
them is an LLM in a tool loop: each tool call costs a full model
generation--call--read cycle, roughly 1--3 \emph{seconds} of model latency.
A question answered with five sequential tool calls costs $5\times$ that ---
10--15 seconds --- regardless of how fast the database is. So:

\begin{center}
\fbox{\parbox{0.92\linewidth}{\textbf{The unit of optimization is the number
of LLM$\leftrightarrow$tool round trips, not per-query database time.} A
primitive that returns a complete, denormalized answer in one call beats five
elegant fine-grained calls, always.}}
\end{center}

This is also the consensus of every deployed system we reviewed in the June
2026 research pass: Microsoft GraphRAG precomputes per-community summary
reports so query time is ``fetch a summary''; LightRAG's headline result is
\emph{one} retrieval call per query versus hundreds for naive graph search;
Apple's Saga serves precomputed entity cards rather than traversing live.
polilabs' versions of this move are the aggregate primitives
(\code{find\_bills\_defining}, \code{count\_bills},
\code{find\_universe\_bills}) and the one-call bill card
(\S\ref{sec:card}).

\section{Outcome derivation as event sourcing}
\label{sec:eventsourcing}

A bill's ``outcome'' (did it pass?) is not a field anywhere in government
data. What exists is an \emph{action trail}: dated strings like ``Passed/agreed
to in House'' or ``Cloture on the motion to proceed not invoked.'' polilabs
derives the outcome by scanning the trail for typed events and keeping the
evidence:

\[
\text{actions} \xrightarrow{\ \text{regex probes}\ }
\text{events} = [(\text{type}, \text{date}, \text{verbatim evidence})]
\xrightarrow{\ \text{precedence rules}\ } \text{outcome enum}
\]

The enum has 13 values: \code{enacted}, \code{vetoed}, \code{failed\_passage},
\code{failed\_cloture}, \code{passed\_house\_only}, \code{passed\_senate\_only},
\code{passed\_both}, \code{died\_after\_passing\_\{house,senate,both\}},
\code{reported\_no\_floor\_vote}, \code{died\_in\_committee}, \code{pending}.
The events list always travels with the label so an agent asked ``why did it
fail?'' quotes dated evidence instead of inventing a narrative.

\textbf{Worked example (the bug that scale found).} H.R.~9495 (118th):

\begin{enumerate}[itemsep=1pt]
  \item 2024-11-12 --- ``Failed of passage \ldots On motion to suspend the
        rules and pass the bill'' $\to$ event \code{failed\_passage}.
        (Suspension requires a 2/3 majority; the bill got 256--145, a
        majority but not 2/3.)
  \item 2024-11-21 --- ``On passage Passed by the Yeas and Nays: 219--184''
        $\to$ event \code{passed\_house}.
  \item Congress closed with no Senate action.
\end{enumerate}

A naive precedence (``any failure beats any passage'') labels this
\code{failed\_passage} --- wrong; the bill passed the House nine days after
the failed shortcut. The fix compares \emph{last-event dates}: a dated
passage supersedes an earlier failed floor vote
(\S\ref{sec:outcomecode}). The correct label is
\code{died\_after\_passing\_house}. None of the 232 curated bills exercised
this path; bill \#46{,}000-ish of the universe did. Scale is a test suite.

\section{The bipartisanship metric, with worked numbers}
\label{sec:bipartisan}

For any roll call with party splits, define for each major party
$p \in \{D, R\}$ the yea-share among that party's yes+no votes,
$s_p = \text{yea}_p / (\text{yea}_p + \text{nay}_p)$, and
\[
\text{bipartisan\_support} \;=\; \min(s_D,\, s_R) \in [0,1].
\]
1.0 means both parties unanimous yes; 0.0 means at least one party unanimous
no. Display labels are thresholds on this number: \code{party\_line} $<0.10
\le$ \code{cross\_party} $< 0.50 \le$ \code{bipartisan} $< 0.90 \le$
\code{near\_unanimous}.

\textbf{Worked example --- the two lives of H.R.~4346.} The same bill number
carried two very different coalitions:

\begin{itemize}[itemsep=2pt]
  \item House roll 239 (2021, as the America COMPETES shell): 215--207.
        Democrats $s_D = 215/215{+}\ldots \approx 1.0$; Republicans
        essentially zero yeas, $s_R \approx 0.005$. $\min \approx 0.005$:
        \code{party\_line}.
  \item House roll 404 (2022-07-28, final CHIPS concurrence): 243--187.
        $s_D = 219/219 = 1.0$; $s_R = 24/(24{+}187) = 24/211 = 0.114$.
        $\min = 0.114$: \code{cross\_party}.
  \item Senate roll 271 (2022-07-27): 64--33, $s_R = 0.347$:
        \code{cross\_party} verging on bipartisan.
\end{itemize}

Two lessons. First, \emph{the metric is per-vote}; a bill-level number must
choose a vote, and polilabs uses the \emph{final passage-class vote}
(\S\ref{sec:finalpassage}) because vehicle bills carry early votes from a
prior life under another name. Second, the labels are display sugar --- the
number is what you filter and sort on, and you should quote the underlying
splits in prose.

Why $\min$ and not, say, the share of the minority party? Because $\min$ is
symmetric (it does not care who holds the majority) and it operationalizes
``did \emph{both} parties support this'' --- the secret-congress question ---
rather than ``did the losing side defect.''

\section{Coverage as a contract (the under/overcoverage primitive)}
\label{sec:coverage}

Agents fail on coverage in two symmetric ways. \emph{Undercoverage error}:
the agent believes the corpus knows less than it does and abstains wrongly.
\emph{Overcoverage error}: the agent believes the corpus knows more than it
does and asserts from a facet that is only partially populated --- the worse
failure, because it reads as hallucination.

The fix is to make coverage a queryable, live-computed table rather than
prose: for each topic and for the universe, for each data \emph{facet}
(actions, cosponsors, votes, outcome, CRS summary, embeddings, full text),
report \code{bills\_with\_facet / bill\_count}. Live-computed means it cannot
drift from the index. The agent's honesty rules become mechanical: before
aggregating over a facet, check that the facet is total; when it is not, say
which slice the answer rests on. This generalizes the schema's original move
for definitions (explicit gap nodes instead of silent misses) to the corpus
itself.

\part{The data: where every fact comes from}

\section{The government data landscape}
\label{sec:sources}

Everything polilabs knows comes from five public, unauthenticated, primary
sources. No aggregator (and no API key) sits in between --- this matters
because the two civic aggregators researchers used to rely on are dead
(ProPublica Congress API archived 2025; GovTrack bulk data ended 2024).

\begin{longtable}{@{}p{3.3cm}p{5.6cm}p{6.0cm}@{}}
\toprule
Source & What it carries & How polilabs uses it \\
\midrule
GovInfo \textbf{BILLS} bulk XML &
Full bill text at each procedural stage (\code{ih} introduced, \code{eh}
engrossed House, \code{enr} enrolled, \ldots), in USLM or pre-USLM XML &
The curated corpora's \file{bill.xml}; parsed into the section tree;
the most-final available version wins \\[3pt]
GovInfo \textbf{BILLSTATUS} bulk XML &
Per-bill status: every title variant, structured sponsor/cosponsors
(party, state, district), the complete dated action trail, recorded-vote
pointers, public-law numbers, policy area, subjects &
The universe layer (one zip per congress$\times$type, 12 zips total);
\file{billstatus.json} enrichment for curated bills; outcome derivation;
the alias table's main source \\[3pt]
GovInfo \textbf{BILLSUM} bulk XML &
CRS-written summaries, revised at each major action; the latest version
describes the bill's final state &
\code{universe\_summaries} table; the FTS \code{summary} column (topical
search over all 46k bills); the card's \code{summary} field \\[3pt]
House Clerk / Senate LIS \textbf{roll-call XML} &
Per-vote: question, result, per-party totals, every member's position
with party and state &
Party splits and \code{bipartisan\_support} on every recorded vote;
member-level positions for passage-class votes on curated bills \\[3pt]
OLRC \textbf{Popular Name Table} &
The official act-nickname dictionary (popular name $\to$ public law) &
926 alias rows for the covered Congresses, source-tagged
\code{olrc\_popular\_names} \\
\bottomrule
\end{longtable}

Update cadence note: BILLSTATUS bulk files are regenerated continuously by
GPO (the directory listing showed same-day timestamps when we pulled);
BILLSUM lags actions by days-to-weeks because humans at CRS write the
summaries. This is why summary coverage is near-total for the closed 117th
(15{,}221 of 15{,}231 bills) but partial for the 118th (9{,}399) and the
in-progress 119th (4{,}256) --- the facet table reports these numbers rather
than hiding them.

\section{Identity: the two bill-ID schemes and the vehicle-bill trap}
\label{sec:identity}

Every bill has two equivalent IDs: the legacy form \code{117-hr-4346}
(SQLite, file paths, tool arguments) and the URN form
\code{bill:us/117/hr/4346} (the graph, where cross-jurisdiction extensions
need a namespace). Section IDs are
\code{\{bill\_id\}::\{xml\_id\}}, where \code{xml\_id} is the USLM
\code{id} attribute --- globally unique within a bill version, which is what
lets graph rows be scoped to a bill by string prefix.

Two identity facts that will bite anyone who does not know them:

\begin{enumerate}[itemsep=2pt]
  \item \textbf{Bill numbers do not persist across Congresses.} ``Laken Riley
        Act'' names 118-hr-7511, 119-hr-29, and 119-s-5; only the last was
        enacted. This is why \code{lookup\_bill} returns \emph{all} matches
        with \code{is\_ambiguous=True} and the agent is instructed to
        disambiguate, never to silently pick.
  \item \textbf{Vehicle bills.} Congress amends a convenient existing bill
        rather than starting fresh, so a bill's roll calls can belong to a
        prior life (H.R.~4346's 2021 votes are America COMPETES votes), and a
        famous vote can live on a \emph{different} number (KOSA's 91--3
        Senate vote was on the S.~2073 vehicle; S.~1409 itself correctly
        shows no floor vote). The data layer handles this by dating every
        vote and using final-passage votes for bill-level metrics; the
        prompt layer instructs the agent to check dates and questions before
        attributing margins.
\end{enumerate}

\section{The two data layers: universe vs.\ curated topics}
\label{sec:twolayers}

This is the central organizing idea after the June 2026 expansion:

\begin{center}
\begin{tabular}{@{}p{3.2cm}p{5.4cm}p{5.4cm}@{}}
\toprule
 & \textbf{Universe layer} & \textbf{Curated topic corpora} \\
\midrule
Population & every law-track bill (hr, s, hjres, sjres) of the 117th--119th:
\textbf{46{,}168 bills} & hand-curated or criteria-gated subsets:
\textbf{232 bills} in 3 topics \\
Carries & status: outcome + evidence, every known name, sponsor party/state,
cosponsor party counts, roll-call party splits, CRS summary, policy area &
everything the universe has \emph{plus} full text: 167k sections, defined
terms, amendment operations, citations, member-level vote positions \\
Source artifact & \file{data/universe/*.jsonl.gz} (committed, $\sim$15\,MB) &
\file{data/corpus/\{topic\}/\{bill\}/} directories (committed) \\
Answers & ``did X pass? was it bipartisan? what is it about? how many\ldots'' &
``what does Sec.~3(b) say? how does it define `foundation model'? what does
it amend?'' \\
\bottomrule
\end{tabular}
\end{center}

The boundary is stated as data (\code{in\_corpus} on every universe row, the
\code{full\_text} facet in the coverage report) and enforced in prompts: a
status-only card tells the agent \emph{not} to call text tools on that bill.
Topics so far: \code{ai\_governance} (191 bills, criteria-gated:
\file{corpus/inclusion\_criteria.md}), \code{redistricting} (12, hand-curated
seed), \code{secret\_congress} (29, hand-curated passage-dynamics exemplars in
six clusters: \file{corpus/secret\_congress\_criteria.md}).

\part{The storage architecture, layer by layer}

\section{Files are the source of truth}
\label{sec:layers}

\begin{center}
\begin{tabular}{@{}llll@{}}
\toprule
Layer & Location & Contents & Disposable? \\
\midrule
0. Cache & \file{data/cache/} & raw zips, BILLSTATUS XML, roll-call XML & yes (gitignored) \\
1. Truth & \file{data/corpus/}, \file{data/universe/} & per-bill dirs; compact jsonl.gz & \textbf{no (committed)} \\
2. Index & \file{data/polilabs.db} & SQLite: relational + FTS + embeddings & yes (rebuilt) \\
3. Graph & \file{data/polilabs.kuzu} & K\`uzu property graph & yes (rebuilt) \\
\bottomrule
\end{tabular}
\end{center}

A curated bill directory contains: \file{bill.xml} (verbatim text),
\file{metadata.json} (identity, title, sponsor, topic, tier, and --- for
secret\_congress --- \code{cluster} and \code{curator\_note}),
\file{billstatus.json} (the keyless enrichment: outcome, events, votes with
members, actions, cosponsors), \file{provenance.json} (source URLs, fetch
times, criteria version). The rule ``never modify \file{metadata.json} after
promotion; enrichment goes in sibling files'' is what made it safe to backfill
all 203 pre-existing bills without touching anything the original pipeline
wrote.

\section{The SQLite index: every table and its job}
\label{sec:sqlite}

Eighteen tables/virtual tables. Grouped by job:

\textbf{Curated-corpus relational core:} \code{bills} (one row per curated
bill; carries outcome, public\_law, bipartisan\_support, cluster,
curator\_note, outcome\_events JSON), \code{sections} (167{,}742 rows; the
parsed tree with \code{parent\_section\_id}, verbatim \code{text\_self} and
recursive \code{text\_full}, canonical citation), \code{bill\_versions},
\code{cosponsors}, \code{actions}, \code{subjects}, \code{citations}
(reserved), \code{corpus\_meta}, \code{source\_freshness}.

\textbf{Votes:} \code{votes} (347 rows; per roll call: chamber, roll number,
question, result, vote type, totals, per-party splits, bipartisan metric,
source URL) and \code{vote\_positions} (20{,}505 rows; member, party, state,
position --- passage-class votes only, which is what makes ``how did the
California delegation vote on CHIPS?'' a one-join query).

\textbf{Universe:} \code{universe\_bills} (46{,}168 rows; the status record),
\code{universe\_votes} (2{,}982 rows; party splits, no members),
\code{universe\_summaries} (28{,}786 rows; latest CRS summary),
\code{bill\_aliases} (97{,}401 rows; \code{alias\_norm} $\to$
\code{bill\_id}, source-tagged).

\textbf{Search:} \code{bills\_fts}, \code{sections\_fts} (corpus FTS;
\code{topic} carried UNINDEXED so query-time filtering is free),
\code{universe\_fts} (title + all aliases + CRS summary --- this is what
makes the whole universe topically searchable), \code{section\_embeddings}
(the dense leg; partial coverage, reported honestly).

The universe row, quoted from \file{index/schema.py}:

\begin{lstlisting}[style=sql]
CREATE TABLE IF NOT EXISTS universe_bills (
    bill_id            TEXT PRIMARY KEY,
    congress           INTEGER NOT NULL,
    bill_type          TEXT NOT NULL,
    bill_number        INTEGER NOT NULL,
    title              TEXT,
    introduced_date    TEXT,
    origin_chamber     TEXT,
    policy_area        TEXT,
    sponsor_name       TEXT,
    sponsor_party      TEXT,
    sponsor_state      TEXT,
    sponsor_bioguide   TEXT,
    cosponsors_d       INTEGER NOT NULL DEFAULT 0,
    cosponsors_r       INTEGER NOT NULL DEFAULT 0,
    cosponsors_i       INTEGER NOT NULL DEFAULT 0,
    cosponsors_total   INTEGER NOT NULL DEFAULT 0,
    latest_action_date TEXT,
    latest_action_text TEXT,
    outcome            TEXT,
    public_law         TEXT,
    outcome_events     TEXT,   -- JSON list, capped
    vote_refs          TEXT,   -- JSON list of recorded-vote pointers
    bipartisan_support REAL,
    bipartisan_label   TEXT,
    in_corpus          INTEGER NOT NULL DEFAULT 0
);
\end{lstlisting}

Note what is \emph{not} here: no text, no sections. A 46k-row table of this
shape is a few tens of megabytes and every query against it is milliseconds.

\section{The K\`uzu graph: nodes, edges, and what is actually populated}
\label{sec:kuzu}

The graph implements the curated corpora's structural layer. Node tables
actually populated today: \code{Bill} (232; now carrying \code{outcome},
\code{public\_law}, \code{bipartisan\_support}, \code{cluster}),
\code{BillVersion} (230), \code{Section} (167{,}742), \code{Sponsor} (747),
\code{StatuteSection} (3{,}188 lazily-merged US Code targets),
\code{DefinedTerm} (3{,}749), \code{AmendmentOperation} (6{,}875),
\code{RollCallVote} (347), plus provenance plumbing. Edge tables populated:
\code{HAS\_VERSION}, \code{HAS\_SECTION}, \code{PARENT\_OF} (166{,}736 ---
the section tree), \code{SPONSORED\_BY}, \code{COSPONSORED\_BY},
\code{CITES\_EXTERNAL} (9{,}412), \code{DEFINES} (3{,}689),
\code{BY\_REFERENCE} (387), \code{AMENDS} (6{,}612), \code{TARGETS}
(6{,}842), \code{HAS\_ROLLCALL} (347).

Declared but still empty (deliberately, with the gaps listed in
\code{corpus\_coverage}): \code{CITES\_INTERNAL} (needs USLM
\code{<ref>} parsing), \code{RESOLVED\_TO}/\code{UnresolvedTermUse}
(definition use-site resolution), \code{REFERRED\_TO} committees.

\section{Killing the 3.5-hour rebuild: the incremental path}
\label{sec:incremental}

The full graph build takes $\sim$3.5\,h on the Mac, dominated by parsing 167k
sections through the citation/definition/amendment extractors. The June
research pass established this is self-imposed, not a K\`uzu limitation:
since v0.4, K\`uzu supports \code{COPY FROM} into \emph{non-empty} tables, the
documented fast path for bulk appends. The target maintenance discipline:

\begin{enumerate}[itemsep=2pt]
  \item New bills $\to$ stage their node/edge rows as Parquet/CSV deltas
        $\to$ \code{COPY FROM} into existing tables. Seconds, not hours.
  \item Property drift on existing nodes (an outcome changes when a pending
        bill passes) $\to$ small Cypher \code{SET} transactions. We already
        run this pattern: a drift-check script compares SQLite
        \code{bills.outcome} against graph properties and patches differences
        --- the date-fix round found zero drifted corpus bills, verified in
        seconds rather than re-parsing for hours.
  \item Full rebuilds only for schema changes.
\end{enumerate}

The universe deliberately lives in SQLite only for now: every universe
primitive is a filtered scan or aggregate (\S\ref{sec:graphvsrel}), so graph
residency would add build cost without adding query power. Universe Bill
nodes go into the graph (via \code{COPY FROM}) when there is a traversal
that needs them --- e.g., same-subject edges across the full universe.

\part{The evolution: how we got here}

\section{The chain at a glance}
\label{sec:chain}

\begin{center}
\begin{tabular}{@{}llp{8.6cm}@{}}
\toprule
Step & When & What it added \\
\midrule
v0.1 schema & spring & 23-type design doc: bitemporal versioning, local
definitions, reified amendments, derivation-typed edges \\
PR1 & spring & bibliographic spine in K\`uzu (Bill/Version/Section/Sponsor) \\
PR2 & spring & \code{CITES\_EXTERNAL} to US Code, statute targets lazily merged \\
PR3 & May & \code{DefinedTerm} + \code{DEFINES} + \code{BY\_REFERENCE};
aggregate definition tools \\
PR4 & May & reified \code{AmendmentOperation} + \code{AMENDS}/\code{TARGETS} \\
Topics (P1--P3) & May & \code{topic} column end-to-end; redistricting seed (12 bills) \\
Round 1 & Jun 9--10 & secret\_congress topic (29 bills, 6 clusters);
\file{billstatus.json} enrichment for all bills; votes + outcomes + facet
coverage; \code{count\_bills}, \code{get\_bill\_votes},
\code{find\_bills\_by\_outcome} \\
Round 2 & Jun 10--12 & the universe layer (46{,}168 bills); alias resolution
(97{,}401 names); one-call bill cards; CRS summaries (28{,}876);
universe-wide FTS; landing-page truthfulness \\
\bottomrule
\end{tabular}
\end{center}

\section{Why each step happened (the motivating failures)}
\label{sec:why}

\textbf{v0.1's bet: definitions are local.} The original schema's deepest
insight, worth restating because everything since builds on its style: there
is no global ``AI system'' node, because two bills' identical surface forms
carry different meanings. Definitions are scoped to their defining section;
cross-bill questions are answered by \emph{aggregating scoped facts}, never by
merging them. Every later layer copies this pattern --- curated cluster tags
sit \emph{beside} mechanical vote records, not instead of them; outcomes carry
their evidence; coverage is per-facet, not a single number.

\textbf{The eval drove the aggregates.} The 20-question agent eval (Opus 4.7:
85\%, Sonnet 4.6: 70\%) failed systematically on one shape: set-valued and
counting questions answered by paginated search + client-side counting.
PR3/PR4 added \code{find\_bills\_defining}/\code{find\_bills\_amending} for
the definitional/amendatory versions of that shape; Round 1's
\code{count\_bills} closed the literal counting hole. The general law:
\emph{when the question shape is known, give the agent a primitive that
returns the complete answer with a coverage note}.

\textbf{The meeting notes drove passage dynamics.} ``Show negative bills that
failed (why it fails?)''; ``a lot of legislation is passed on a bipartisan
vote in a single bill''; ``how much it affected the vote.'' These are not text
questions, and the representation had no nouns for them --- no outcome field,
no votes, action trails on only $\sim$30\% of bills. Round 1 added the nouns
mechanically (BILLSTATUS + roll-call XML are keyless), then curated 29
exemplars into the clusters the notes named (quiet bipartisan laws, bipartisan
bills that died, omnibus vehicles, party-line contrasts).

\textbf{Scale made the thesis testable, and the lack of scale was the
critique.} With only 29 curated exemplars, ``bipartisan laws pass quietly''
is an assertion. With the universe layer it is a query, and the answer is
striking: of the 639 laws enacted by the 117th--118th, \textbf{361 (57\%)
passed with no roll call at all} (voice vote or unanimous consent --- the
quietest possible path), 244 passed bipartisan-to-near-unanimous, and only
\textbf{13 (2\%) were party-line} --- exactly the reconciliation bills, CRA
disapproval resolutions, and appropriations the thesis predicts. The curated
clusters remain as \emph{interpretation}; the universe is the
\emph{denominator}.

\textbf{The research pass drove the interface restructure.} The June review
of deployed KG systems (GraphRAG, LightRAG, Saga, Wikidata/UMLS alias
practice, the one published legislative KG, K\`uzu's own docs) returned one
contradiction with our design: fine-grained primitives put graph traversal
\emph{inside} the LLM round-trip loop (\S\ref{sec:latency}). The responses:
alias-table name resolution (the literature's unanimous first step),
the one-call card (Saga's entity-card pattern at our scale), free-text
search over the universe (so topical discovery does not require knowing
which curated topic to search), and the routing rules in the system prompt.

\part{The query surface: all 18 tools}

\section{The tool table}
\label{sec:tools}

Grouped by the question shape they serve. ``Layer'' says which storage
answers it.

\begin{longtable}{@{}p{3.6cm}p{1.6cm}p{9.2cm}@{}}
\toprule
Tool & Layer & Job (and when the agent should reach for it) \\
\midrule
\multicolumn{3}{@{}l}{\emph{Entry points}} \\
\code{lookup\_bill} & SQLite & name/nickname/id $\to$ canonical bill\_id(s);
ALWAYS first when the user names a bill \\
\code{get\_bill\_card} & SQLite{+}K\`uzu & one-call card: names, sponsor,
outcome+evidence, votes, CRS summary, structure counts, in\_corpus flag \\
\code{search\_corpus} & SQLite & hybrid BM25+dense over curated full text,
topic-scoped \\
\code{corpus\_coverage} & SQLite & the coverage contract: topics, facets,
universe split, known gaps \\
\midrule
\multicolumn{3}{@{}l}{\emph{Set-valued / counting (one call = complete answer)}} \\
\code{find\_universe\_bills} & SQLite & filters + free-text query over all
46k bills; the denominator tool \\
\code{count\_bills} & SQLite & exact counts over curated bills, group-by
topic/congress/outcome/\ldots \\
\code{find\_bills\_by\_outcome} & SQLite & curated bills by outcome /
cluster / bipartisanship band \\
\code{find\_bills\_defining} & K\`uzu & every bill defining a term
(+ by-reference target filters) \\
\code{find\_bills\_amending} & K\`uzu & per-bill rollup of who amends a USC
section \\
\code{find\_definitions\_of} & K\`uzu & every bill's verbatim definition of a
term, side by side \\
\midrule
\multicolumn{3}{@{}l}{\emph{Single-entity drill-down (curated bills only)}} \\
\code{get\_bill} & SQLite & metadata + section table of contents (no bodies) \\
\code{get\_section} & SQLite & verbatim text + canonical citation \\
\code{get\_bill\_votes} & SQLite & full vote record + outcome + events for
one bill \\
\code{get\_defined\_terms} & K\`uzu & every term one bill defines \\
\code{get\_amendments} & K\`uzu & every amendment operation one bill issues \\
\code{get\_amendments\_targeting} & K\`uzu & operation-level detail on a USC
section \\
\code{get\_citation\_graph} & K\`uzu & typed citation neighborhood of a section \\
\code{resolve\_citation} & SQLite & prose citation $\to$ section\_id \\
\bottomrule
\end{longtable}

\section{Routing: how the agent picks its first tool}
\label{sec:routing}

The system prompt encodes a decision tree, because the wrong first tool costs
round trips that no database speed can recover:

\begin{itemize}[itemsep=2pt]
  \item User \textbf{names a bill} (``the CHIPS Act,'' ``H.R.~4346'') $\to$
        \code{lookup\_bill}, then \code{get\_bill\_card} on the resolved id.
        The card usually ends the question in two calls and reports whether
        full-text tools will work before any are tried.
  \item User asks a \textbf{topical/text} question $\to$
        \code{search\_corpus} on the right topic for deep passages, or
        \code{find\_universe\_bills(query=\ldots)} for breadth across all
        46k bills (titles + aliases + CRS summaries).
  \item User asks a \textbf{set or count} question $\to$ the aggregate tools,
        never paginated search.
  \item Ambiguity from \code{lookup\_bill} is surfaced, not swallowed;
        \code{in\_corpus=False} cards forbid text drill-down on that bill.
\end{itemize}

\part{The code, walked}

\section{Primer: the two query languages}
\label{sec:primer}

\textbf{SQLite FTS5.} \code{CREATE VIRTUAL TABLE x USING fts5(col, \ldots,
tokenize='porter unicode61')} builds the inverted index;
\code{WHERE x MATCH 'fentanyl scheduling'} returns rows containing both
stems; \code{ORDER BY rank} sorts by BM25. A column declared
\code{UNINDEXED} is stored but not searchable --- we use that for
\code{bill\_id}/\code{topic} so they can be filtered/joined without
polluting relevance.

\textbf{K\`uzu Cypher.} \code{MATCH (b:Bill \{bill\_id: \$bid\})-[:HAS\_ROLLCALL]->(v:RollCallVote)
RETURN v} binds node patterns by label and property, follows typed edges,
returns rows. \code{STARTS WITH} on a primary-key-ish string is how we scope
section-keyed rows to a bill. Parameters (\code{\$bid}) are passed as a dict
--- never interpolate strings into Cypher.

\section{Outcome derivation (\file{ingest/billstatus.py})}
\label{sec:outcomecode}

The probes are compiled regexes over action text:

\begin{lstlisting}
_PASSED_HOUSE  = re.compile(r"passed[ /]agreed to in house|^passed house|on passage passed", re.IGNORECASE)
_PASSED_SENATE = re.compile(r"passed[ /]agreed to in senate|^passed senate", re.IGNORECASE)
_FAILED_PASSAGE = re.compile(r"failed of passage|on passage failed|failed passage", re.IGNORECASE)
_FAILED_CLOTURE = re.compile(r"cloture(?: on the )?.*not invoked|cloture motion .* rejected", re.IGNORECASE)
_VETOED   = re.compile(r"vetoed by president", re.IGNORECASE)
_REPORTED = re.compile(r"reported (?:to|by)|placed on .* calendar", re.IGNORECASE)
\end{lstlisting}

Each match appends a dated, deduplicated event (BILLSTATUS repeats the same
event from multiple source systems; we keep one per (type, day)). Roll-call
results add failure evidence that action strings sometimes lack. Then the
date-aware precedence:

\begin{lstlisting}
# A failed floor vote is terminal only if nothing passed AFTER it.
def _last_date(*names: str) -> str:
    return max((e["date"] or "" for e in events if e["event"] in names), default="")

last_pass = _last_date("passed_house", "passed_senate")
failed_passage_final = failed_passage and _last_date("failed_passage") >= last_pass
failed_cloture_final = failed_cloture and _last_date("failed_cloture") >= last_pass

if laws:                outcome = "enacted"
elif vetoed:            outcome = "vetoed"
elif failed_passage_final and congress_closed:  outcome = "failed_passage"
elif failed_cloture_final and congress_closed:  outcome = "failed_cloture"
elif passed_house and passed_senate:
    outcome = "passed_both" if not congress_closed else "died_after_passing_both"
...
\end{lstlisting}

Line by line: \code{laws} (the BILLSTATUS \code{<laws>} element) trumps
everything --- enactment is the one fact the government states directly.
A veto outranks floor failures. A failure is ``final'' only if no chamber
passage postdates it (the H.R.~9495 lesson, \S\ref{sec:eventsourcing}).
Chamber-passage flags then split by whether the Congress has adjourned:
the same facts mean \code{passed\_house\_only} (alive) or
\code{died\_after\_passing\_house} (dead) depending on
\code{congress\_closed} --- a frozen set of past Congresses that must be
updated when the 119th ends (this is the one piece of calendar state in the
file).

\section{The alias contract (\file{index/build\_universe\_tables.py} +
\file{api/\_impl.py})}
\label{sec:lookupcode}

Build-time normalization --- the contract both sides must share:

\begin{lstlisting}
def normalize_alias(s: str) -> str:
    s = s.lower()
    s = _PUNCT.sub(" ", s)        # punctuation -> spaces
    s = _WS.sub(" ", s).strip()   # collapse whitespace
    if s.startswith("the "):
        s = s[4:]
    s = _YEAR_SUFFIX.sub("", s)   # drop trailing 'of 2022'
    return s
\end{lstlisting}

Worked example: \code{"The CHIPS and Science Act of 2022"} $\to$ lowercase
$\to$ punctuation/whitespace pass $\to$ \code{"the chips and science act of
2022"} $\to$ drop leading article $\to$ \code{"chips and science act of
2022"} $\to$ drop year suffix $\to$ \code{"chips and science act"}. The same
function runs at query time (imported from the same module --- one
definition, two call sites), so ``the Chips \& Science Act'' from a user
normalizes to an identical key. The year-drop is what merges ``CHIPS Act of
2022'' with ``CHIPS Act''; the table keeps the original surface form so
results can show what actually matched.

Query-time ladder in \code{lookup\_bill}: (1) bill-id forms via the two
ID regexes; (2) exact \code{alias\_norm} match (the dictionary step that
resolves the majority of mentions); (3) FTS5 over
title+aliases(+summaries) as the fuzzy fallback. Each rung is tried only if
the previous returned nothing; measured cost is 0.1--20\,ms.

\section{The bill-level bipartisan number (\file{index/build.py})}
\label{sec:finalpassage}

\begin{lstlisting}
def _final_passage_bipartisan(bs: dict | None) -> float | None:
    """bills.bipartisan_support = the metric on the LAST passage-class
    roll call (chronologically). Vehicle bills carry early votes from a
    prior life under another name; the final vote is the one that
    reflects the enacted coalition."""
    if not bs:
        return None
    passage = [v for v in bs.get("votes", [])
               if v.get("vote_type") == "passage"
               and v.get("bipartisan_support") is not None]
    if not passage:
        return None
    passage.sort(key=lambda v: v.get("date") or "")
    return passage[-1]["bipartisan_support"]
\end{lstlisting}

Three deliberate exclusions, each an honesty decision: amendment and
procedural votes never represent the coalition for the bill (motions to
recommit are designed to split parties); votes without party decomposition
(voice votes have no roll at all) yield \code{None}, not 0 --- \code{None}
means ``no partisan-decomposable evidence,'' which is itself informative
(57\% of enacted laws); and \code{None} rows are excluded by band filters so
``bipartisan\_support $\ge$ 0.5'' never silently includes unknowns.

\section{The card assembly (\file{api/\_impl.py::get\_bill\_card})}
\label{sec:card}

The card is assembled at read time from five cheap queries (universe row,
corpus row, aliases, votes, CRS summary) plus two optional graph counts:

\begin{lstlisting}
# Graph section ids are '{urn_bill_id}::{xml_id}', so a prefix
# match scopes DefinedTerm / AmendmentOperation rows to this
# bill. Failures degrade to None (= unknown), never to 0.
urn_prefix = _to_urn_bill_id(bill_id) + "::"
for query, which in (
    ("MATCH (d:DefinedTerm) WHERE d.defining_section_id STARTS WITH $p RETURN COUNT(d)", "terms"),
    ("MATCH (a:AmendmentOperation) WHERE a.source_section_id STARTS WITH $p RETURN COUNT(a)", "amendments"),
):
    ...
\end{lstlisting}

Two details worth internalizing. First, the \code{None}-vs-0 discipline
again: if the graph is unavailable, counts are \code{None} (``unknown''),
because returning 0 would teach the agent the bill defines nothing. Second,
why \emph{compute-on-read} rather than precomputing card blobs (as Saga
does): every component query is single-digit milliseconds at our scale, so
precomputation would add a staleness liability and buy no agent-visible
latency --- the round trips were the cost, and one call is one call either
way. If per-card cost ever grows (e.g., neighbor summaries), the switch to a
precomputed \code{bill\_cards} table is mechanical.

\section{The universe build (\file{scripts/build\_universe.py})}
\label{sec:universecode}

The pipeline that makes 46k bills cost 19 seconds: GPO publishes one zip per
(congress, bill type); the build streams each zip member through the same
\code{parse\_billstatus} used everywhere else, compacts to the record shape
(cosponsor lists become party \emph{counts}; events capped at 12; subjects at
12), skips the few auxiliary zip members with no bill identity (they would
collapse into a garbage \code{117--0} row --- found the hard way), sorts for
deterministic diffs, and writes gzipped JSONL. Roll calls are fetched once
per \emph{unique vote URL} (2{,}982 across all three Congresses) with a disk
cache, then joined back to every bill that references the URL. Summaries
(BILLSUM) get the same treatment with a latest-version pick. Everything is
re-runnable and idempotent; the zips are gitignored cache, the jsonl.gz
artifacts are committed truth.

\part{Cheat sheet: your questions, answered}

\section{Index of questions}
\label{sec:cheat}

\begin{description}[itemsep=3pt]
  \item[``You built this off an older branch?''] No --- verified zero commits
    behind \code{origin/master} at both build rounds; the branch is
    master + the two overnight commits. It \emph{looked} old because rounds
    1--2 were backend/data work with no visible pixels, and the landing page
    showed stale hardcoded counts --- fixed by making the stats pipeline live
    (\S\ref{sec:coverage}, landing section of the PR).
  \item[``If there's 46{,}000 bills it needs to say 46{,}000.''] Agreed and
    enforced: the landing fixture is regenerated from the index
    (46{,}168 tracked / 739 enacted / 3{,}329 roll calls / 232 full-text),
    and the numbers ship with a cache-bust version so deployed browsers
    cannot show stale ones. \S\ref{sec:twolayers} for what each number means.
  \item[``Where are we pulling the bills from?''] Five keyless primary
    sources, now also stated on the landing page: \S\ref{sec:sources}.
  \item[``Danny's additional statistics --- legislator notes / LA notes?''] I
    read this as the CRS-written bill summaries in GovInfo's BILLSUM
    collection (28{,}876 ingested; on every card and searchable:
    \S\ref{sec:sources}). If ``LA notes'' meant something else --- e.g.,
    committee-report text or CBO cost estimates --- both exist as GovInfo
    collections and would follow the same ingest pattern; flag it and it's a
    day's work, not a redesign.
  \item[``How well the bill was passed?''] That is \code{bipartisan\_support}
    plus the per-vote party splits: \S\ref{sec:bipartisan}. Margins, party
    splits, and the no-roll-call (voice/UC) case are all distinguished.
  \item[``Impact on sponsorship / donors?''] Donor \emph{behavior} data (FEC)
    is not ingested --- it is a different identity universe (committees,
    filings) and a real research-design question, not an ingest script. The
    donor \emph{topic} is represented (DISCLOSE Act, the platform-lobbying
    deaths of AICO/JCPA). \S\ref{sec:why}.
  \item[``Can you get vote shares?''] Electoral vote shares post-vote: out of
    scope tonight, agreed it's the biggest ask. The path exists --- member
    positions are keyed by bioguide ID, MIT Election Lab has results ---
    but linking them honestly (redistricting, retirements, identification)
    needs a design discussion first.
  \item[``Why did bill X fail?''] \code{get\_bill\_card} or
    \code{get\_bill\_votes}: outcome + dated evidence events
    (\S\ref{sec:eventsourcing}); for sets, \code{find\_bills\_by\_outcome} /
    \code{find\_universe\_bills}.
  \item[``Was the CHIPS Act bipartisan?''] Depends which vote --- the worked
    example in \S\ref{sec:bipartisan} is exactly this bill, and the
    vehicle-bill trap (\S\ref{sec:identity}) is why the question needs the
    per-vote answer.
  \item[``What is `secret congress' in our data?''] A curated 29-bill topic
    (six clusters) \emph{plus} a now-testable universe-level claim: 57\% of
    117th--118th laws passed with no roll call, 2\% party-line.
    \S\ref{sec:why}.
  \item[``Why is the Kuzu rebuild 3.5 hours and do we care?''] Extractor
    parsing, not graph insertion; and we stop caring once appends go through
    \code{COPY FROM} deltas and property drift goes through Cypher
    \code{SET}: \S\ref{sec:incremental}.
  \item[``What keeps the agent from making things up about coverage?''] The
    facet-coverage contract in \code{corpus\_coverage} plus prompts that
    treat \code{None} as ``unknown'' and status-only cards as text-free:
    \S\ref{sec:coverage}, \S\ref{sec:card}.
  \item[``What's still missing?''] Embeddings (handed off), CITES\_INTERNAL
    and definition use-sites (declared, empty, listed in known gaps),
    committee/subject edges, universe nodes in the graph, measured media
    salience for the secret-congress clusters, FEC/electoral linkage.
    Each has a stated path; none blocks the current surface.
\end{description}

\end{document}
